\esSezione{Multiplexer 16:1}\label{sec:esercizio1-1}

\begin{traccia}
Progettare, implementare in VHDL e testare mediante simulazione un multiplexer 16:1 a un bit, con sedici ingressi dati e quattro ingressi di selezione, realizzato in modo strutturale a partire da multiplexer 4:1.
\end{traccia}

\progetto

Il multiplexer 16:1 è costruito su due livelli di multiplexer 4:1. Al primo livello ci sono quattro componenti: ognuno riceve un gruppo di quattro ingressi consecutivi e i due bit meno significativi della selezione. Le loro uscite alimentano un quinto multiplexer 4:1, pilotato dai due bit più significativi della selezione, e l'uscita di quest'ultimo è quella del top module.

% FIXME: schema a blocchi non presente tra le immagini del progetto.

\implementazione

Partiamo dal componente elementare, l'entità \texttt{mux\_4\_1}, che ha come porte \texttt{a}, vettore di quattro ingressi dati indicizzato da 0 a 3, \texttt{s}, selezione a due bit, e \texttt{y}, uscita a un bit. L'architettura \texttt{dataflow} usa un'assegnazione \texttt{with select}: a ogni valore di \texttt{s} corrisponde un ingresso, mentre per i valori non validi l'uscita assume \texttt{'-'}. Il codice è il seguente.

\vhdlfile{esercizio1-1/mux_4_1.vhd}{Implementazione MUX 4:1}{lst:esercizio1-1:mux-4-1}

Definito tale componente, vediamo il top module \texttt{mux\_16\_1}, che ha come porte il vettore \texttt{a} di sedici ingressi (indici da 0 a 15), la selezione \texttt{s} a quattro bit e l'uscita \texttt{y}. L'architettura è \texttt{structural}: il segnale interno \texttt{u}, di quattro bit, raccoglie le uscite dei componenti \texttt{M0}, \texttt{M1}, \texttt{M2} ed \texttt{M3}, istanziati con port mapping sui gruppi \texttt{a(0 to 3)}, \texttt{a(4 to 7)}, \texttt{a(8 to 11)} e \texttt{a(12 to 15)} e con \texttt{s(1 downto 0)} come selezione. Il componente \texttt{M4} riceve \texttt{u} come vettore di ingressi, ha \texttt{s(3 downto 2)} come selezione e produce \texttt{y}.

\vhdlfile{esercizio1-1/mux_16_1.vhd}{Implementazione MUX 16:1}{lst:esercizio1-1:mux-16-1}

L'entità \texttt{mux\_4\_1} non verrà ripresentata quando sarà riutilizzata in altri esercizi.

\simulazione

% TODO: testbench mancante (vivado/sim/mux_16_1_tb.vhd, la cartella vivado/sim non esiste nel progetto).
\todo{Testbench \texttt{mux\_16\_1\_tb.vhd} non disponibile: listato e commento della sua struttura da aggiungere.}

In mancanza del testbench si commenta la sola forma d'onda ottenuta. Nella figura~\ref{fig:esercizio1-1:simulazione} l'ingresso, indicato nella finestra come \texttt{i}, resta costante al valore 0x2D4A ($0010\,1101\,0100\,1010_2$), mentre la selezione \texttt{s} assume in successione i valori da 0x0 a 0xE, con un passo di \SI{10}{ns}, fino a 0xF.

\begin{figure}[H]
  \centering
  \includegraphics[width=0.95\linewidth]{figure/esercizio1-1/simulazione}
  \caption{Simulazione MUX 16:1}
  \label{fig:esercizio1-1:simulazione}
\end{figure}

Poiché l'indice 0 di \texttt{a} corrisponde al bit più a sinistra, con \texttt{s} pari a 0x2 ($0010_2$) l'uscita è \texttt{a(2)}, ossia \texttt{'1'}, e infatti \texttt{y} è alto per tutto l'intervallo in cui la selezione vale 2. Con \texttt{s} pari a 0xF ($1111_2$) viene selezionato \texttt{a(15)}, uguale a \texttt{'0'}, e l'uscita osservata è \texttt{'0'}.
